\section{Καμπύλες Πλήρωσης Χώρου Hilbert}
\label{sec:bg_hilbert}

\subsection*{Τυπικός Ορισμός και Αναδρομική Κατασκευή}

Οι \textit{Καμπύλες Πλήρωσης Χώρου} (Space-Filling Curves --- SFCs) παρέχουν συνεχή, επί (surjective) απεικόνιση από τον μοναδιαίο υπερκύβο $[0,1]^d$ στο μοναδιαίο διάστημα $[0,1]$~\cite{sagan1994space}. Η \textit{καμπύλη Hilbert} τάξης $p$ στις $d$ διαστάσεις ορίζεται αναδρομικά ως η οριακή συνάρτηση μιας σειράς λεπτομερέστερων προσεγγίσεων $\mathcal{H}_p^{(d)} : [0, 2^{pd})\to [0,1]^d$. Στην $p$-οστή τάξη, ο χώρος $[0,1]^d$ υποδιαιρείται σε $2^{pd}$ ίσα υπερκυβικά κελιά (hypercubes), τα οποία διατάσσονται σε μια συνεχή αλυσίδα U-σχήματος· κάθε κελί αντιστοιχεί σε έναν μοναδικό ακέραιο δείκτη $h \in \{0, 1, \ldots, 2^{pd}-1\}$. Τυπικά:
\begin{equation}
    \mathcal{H}_p^{(d)} : \{0,\ldots,2^{pd}-1\} \to \{0,\ldots,2^p-1\}^d,
\end{equation}
και η απεικόνιση $\mathcal{H}_p^{(d)}$ είναι βιεκτική (bijection) μεταξύ ακεραίων δεικτών και κελιών πλέγματος.

\subsection*{Θεώρημα Διατήρησης Τοπικότητας}

Η θεμελιώδης ιδιότητα της καμπύλης Hilbert είναι η \textit{διατήρηση τοπικότητας} (locality preservation): σημεία που βρίσκονται κοντά στον πολυδιάστατο χώρο $[0,1]^d$ αντιστοιχίζονται, κατά το δυνατόν, σε κοντινούς δείκτες $h$ στη μονοδιάστατη ακολουθία~\cite{lawder2000spacefilling, sagan1994space}. Πιο συγκεκριμένα, για δύο σημεία $\bm{x}, \bm{x}' \in [0,1]^d$ με Ευκλείδεια απόσταση $\|\bm{x} - \bm{x}'\|_2 \leq \epsilon$, ισχύει:
\begin{equation}
    \bigl|\mathcal{H}_p(\bm{x}) - \mathcal{H}_p(\bm{x}')\bigr| = O\!\left(\epsilon^{1/d} \cdot 2^{pd}\right).
\end{equation}
Αντιστρόφως, γειτονικοί δείκτες $h, h'$ αντιστοιχούν σε γειτονικά κελιά στον πολυδιάστατο χώρο. Η ιδιότητα αυτή δεν ισχύει ποτέ τέλεια (καμία συνεχής SFC δεν μπορεί να διατηρεί τέλεια την τοπικότητα παντού, λόγω τοπολογικών περιορισμών), αλλά η καμπύλη Hilbert αποδεδειγμένα υπερτερεί της \textit{καμπύλης Z-order} (Morton order) σε ό,τι αφορά τη μέση αύξηση της απόστασης μεταξύ γειτονικών κελιών: ενώ η Z-order εισάγει μεγάλα «άλματα» στα όρια των κυρίων υπομπλοκ, η Hilbert τα αποφεύγει χάρη στην U-συμμετρία της ανακατασκευής.

\subsection*{Αλγόριθμος Bit-Manipulation του Skilling}

Η αποδοτική υπολογιστική υλοποίηση της απεικόνισης Hilbert σε αυθαίρετες διαστάσεις $d$ και τάξη $p$ επιτεύχθηκε από τον Skilling~\cite{skilling2004hilbert} μέσω άμεσης χειραγώγησης bits (bit-manipulation algorithm). Αντί να αποθηκεύει ρητά τον πίνακα μετατροπής, ο αλγόριθμος εκτελεί μια σειρά από XOR, αριστερές/δεξιές ολισθήσεις και ανταλλαγές bits για να μετατρέψει τις $d \cdot p$ bits συντεταγμένης σε έναν $d \cdot p$-bit ακέραιο δείκτη Hilbert σε $O(d \cdot p)$ χρόνο. Η γραμμική πολυπλοκότητα τον καθιστά κατάλληλο για χρήση σε πραγματικό χρόνο ακόμη και σε υψηλές διαστάσεις ($d \leq 16$) και λεπτομέρεια ($p \leq 32$).

\subsection*{Multi-Probe Hilbert και Ασυνέχειες Ορίων}

Ένας σοβαρός περιορισμός κάθε SFC είναι οι \textit{ασυνέχειες στα όρια} (boundary discontinuities): σημεία εκατέρωθεν ενός ορίου υποδιαίρεσης μπορεί να λαμβάνουν δείκτες που απέχουν $O(2^{pd})$ μεταξύ τους, παρά την πολύ μικρή Ευκλείδεια απόστασή τους. Το φαινόμενο επιδεινώνεται σε υψηλές διαστάσεις, όπου αυξάνεται ο αριθμός των υπερ-επιφανειών ορίων. Για να αντιμετωπιστεί αυτό, η τεχνική \textit{Multi-Probe}~\cite{lv2007multiprobe} εφαρμόζει πολλαπλές μετατοπισμένες και περιστραμμένες εκδόσεις της καμπύλης: ένα σημείο $\bm{x}$ αποκτά $P$ δείκτες $h_1, h_2, \ldots, h_P$ από $P$ διαφορετικές καμπύλες, μεγιστοποιώντας την πιθανότητα ο πραγματικός γείτονας να συμπίπτει με τον υποψήφιο γείτονα σε τουλάχιστον μία καμπύλη.

\subsection*{Μετρική Ποιότητας: Συντελεστής Συσχέτισης Αποστάσεων}

Η ποιότητα της διατήρησης τοπικότητας μετράται με τον \textit{συντελεστή συσχέτισης Pearson} $r$ μεταξύ της Ευκλείδειας απόστασης $\|\bm{x}_i - \bm{x}_j\|_2$ και της διαφοράς δεικτών $|h_i - h_j|$ σε τυχαίο δείγμα ζευγών σημείων:
\begin{equation}
    r = \frac{\sum_{(i,j)}(d^{\text{Euc}}_{ij} - \bar{d}^{\text{Euc}})(d^{\mathcal{H}}_{ij} - \bar{d}^{\mathcal{H}})}{\sqrt{\sum_{(i,j)}(d^{\text{Euc}}_{ij} - \bar{d}^{\text{Euc}})^2\cdot\sum_{(i,j)}(d^{\mathcal{H}}_{ij} - \bar{d}^{\mathcal{H}})^2}}.
\end{equation}
Τιμές $r \to 1$ υποδηλώνουν εξαιρετική διατήρηση τοπικότητας. Εμπειρικές μετρήσεις επιβεβαιώνουν ότι η Multi-Probe Hilbert επιτυγχάνει $r > 0.9$ σε διαστάσεις $d \leq 8$, ενώ η καμπύλη Z-order παρουσιάζει $r$ χαμηλότερο κατά $5{-}15\%$ λόγω των αναφερθέντων ασυνεχειών ορίων.
